\chapter{Glossary of Key Terms}
\addcontentsline{toc}{chapter}{Appendix A: Glossary of Key Terms}

\begin{description}[style=nextline,leftmargin=2em,font=\bfseries\sffamily]
\item[BOM (Bill of Materials)] Complete list of every component in the design with descriptions, part numbers, vendors, costs, quantities, and procurement status. The financial backbone of the build phase.
\item[Build Track] Project track requiring a physical prototype that demonstrates functional performance.
\item[CDR (Critical Design Review)] Design gate answering ``Will this design work?'' Occurs in SDII Sprint~8. Design freezes after CDR.
\item[CF (Course Faculty)] The faculty member(s) responsible for the capstone course. Sets policies, resolves escalated issues, and assigns final grades.
\item[CLO (Course Learning Outcome)] A specific, measurable learning objective that the course is designed to achieve. Capstone has 17 CLOs.
\item[CLT (Capstone Leadership Team)] Program leadership that approves Go/No-Go decisions for travel and resolves program-level issues.
\item[Competition Track] Project track organized around an external competition (SAE, ASCE, AIAA, etc.) with rules functioning as requirements.
\item[ConOps (Concept of Operations)] Narrative description of how the system will be used in its operating environment.
\item[CP (Capstone Purchasing)] The administrative team that processes all purchase orders, manages vendor relationships, and handles reimbursements. Office: Labriola~204.
\item[DFE (Design for Environment)] Sustainability assessment framework for product design, evaluating lifecycle environmental impacts.
\item[Digital Design Archive] Complete, organized submission of all project files at Final Checkout.
\item[Engineering Calculation Package] Collection of formal, documented analyses (hand calculations, FEA, CFD, circuit analysis) supporting the CDR design.
\item[ENVISION] Sustainability assessment framework for infrastructure projects.
\item[ESJ (Engineering for Social Justice)] Framework for assessing equity, access, and social impact in engineering design.
\item[FDVSE (Feasibility, Desirability, Viability, Sustainability, Ethical Responsibility)] Five-dimension value assessment framework for evaluating whether a product is worth building and should be built. A product must satisfy all five dimensions to create genuine value. Introduced in Chapter~5, Section~\ref{sec:fdvse}.
\item[FDR (Final Design Review)] Design gate answering ``How well did it work?'' Occurs in SDII Sprint~13--14.
\item[FE Exam (Fundamentals of Engineering)] A 5-hour-and-20-minute, computer-based exam administered by NCEES. Passing the FE is the first step toward PE licensure and grants the title Engineer Intern (EI).
\item[Feedback Loop] Peer evaluation tool used throughout capstone. Results scale individual grades.
\item[FMEA (Failure Mode and Effects Analysis)] Systematic technique for identifying potential failures and prioritizing mitigation by Risk Priority Number ($\text{RPN} = S \times O \times D$).
\item[FRL (Filter-Regulator-Lubricator)] Air preparation unit for pneumatic systems.
\item[ICD (Interface Control Document)] Specification defining what crosses a subsystem boundary: signals, power, data, format, and constraints.
\item[InnoHub] Labriola Innovation Hub (the primary fabrication facility for capstone teams, including Metal Shop, Wood Shop, Electronics Shop, Composites Shop, and Makerspace).
\item[LEED] Leadership in Energy and Environmental Design (LEED), a sustainability assessment framework for building-related projects.
\item[MoSCoW] Requirements prioritization method: Must Have, Should Have, Could Have, Won't Have (this time).
\item[MVP (Minimum Viable Product)] The smallest set of features that delivers value to the client.
\item[NCEES (National Council of Examiners for Engineering and Surveying)] Organization that develops and administers the FE and PE examinations and publishes the Model Rules of Professional Conduct.
\item[O\&M Manual] Operations and Maintenance manual delivered at FDR, enabling the client to operate, maintain, and troubleshoot the system.
\item[PA (Project Advisor)] The faculty member or industry professional who mentors your team throughout both semesters. Primary point of contact for project-related concerns.
\item[Paper Track] Project track requiring a comprehensive analytical deliverable (simulation, model, technical report) rather than a physical prototype.
\item[PDR (Preliminary Design Review)] Design gate answering ``Should this design work?'' Occurs in SDI Sprint~5. SDRD is baselined after PDR.
\item[PE (Professional Engineer)] Licensed engineering credential requiring an accredited degree, FE exam, 4+ years of progressive engineering experience under a licensed PE, and passage of the PE exam. See Chapter~22.
\item[PIP (Performance Improvement Plan)] Formal process for addressing team member non-participation, with a two-week improvement timeline.
\item[PoC (Proof of Concept)] Demonstration that retires the project's highest technical risk. Presented at PDR.
\item[Product Backlog] Master list of all work the project requires, prioritized by the client.
\item[Scrum Master] Team member who facilitates the sprint process, tracks tasks, runs stand-ups, and removes impediments.
\item[SDI / SDII] Senior Design I (EDNS~491) and Senior Design II (EDNS~492).
\item[SDRD (System Design Requirements Document)] The formal requirements document containing all SHALL/SHOULD statements with IDs, traceability, TAID assignments, and MoSCoW priorities.
\item[SMART Goals] Specific, Measurable, Achievable, Relevant, Time-bound goals set at the start of each semester.
\item[SOW (Statement of Work)] Foundational agreement defining scope, deliverables, schedule, budget, roles, assumptions, and risks.
\item[Sprint Backlog] Subset of Product Backlog items selected for the current sprint.
\item[Sprint Retrospective] End-of-sprint team reflection: what went well, what to improve, what actions to take.
\item[Sprint Review] End-of-sprint demonstration of completed work to the PA and client.
\item[STAR Method] Behavioral interview framework: Situation, Task, Action, Result.
\item[TAID] Verification method categories: Test, Analysis, Inspection, or Demonstration.
\item[Team Contract] Document expressing minimum commitments (attendance, deadlines, contribution, behavior) that all team members agree to uphold.
\item[VCRM (Verification Cross-Reference Matrix)] Table mapping every requirement to its verification method, planned phase, result, and disposition.
\item[WBS (Work Breakdown Structure)] Hierarchical decomposition of the project into progressively smaller tasks.
\end{description}


